\subsection{Increasing a corner by one arrow}
\begin{lemma}[Proved: identity padding preserves the classifying space]
\label{CP2:A-identity-padding}
Let $C$ be a nonempty finite category and $x$ one of its objects.
There is a category $C^+$ with the same objects and support, one
additional endomorphism at $x$, and $BC^+\simeq BC$. Thus the
existence of a model with $n$ arrows and $k$ objects implies
existence with every larger arrow count and the same $k$.
The assertion concerns support; it does not preserve acyclicity
of the category when that term requires trivial endomorphisms.
\end{lemma}
\begin{proof}
In $C^+$ retain every old arrow and old composite, but declare
the old $1_x$ to be an additional idempotent endomorphism and
adjoin a new identity $1_x^+$. The new arrow is a two-sided
identity for every arrow incident with $x$. Associativity
follows from that of $C$: a triple using a new identity reduces
to an old composition, and all other triples are unchanged.
There is exactly one new arrow, only the $x$ corner changes,
and no hom-set becomes empty or nonempty.

Here is a comparison using actual functors. Adjoin to $C$ a
new object $x'$ and form $D$ as follows. Put $r(x')=x$ and
$r(y)=y$ on old objects. For every ordered pair $(a,b)$ except
$(x',x')$, put $D(a,b)=C(r(a),r(b))$, with distinct typed copies
when a new object occurs. Put
$D(x',x')=C(x,x)\sqcup\{1_{x'}\}$. Old-labelled products use
the composition of $C$, and the new $1_{x'}$ is an identity.
This again defines a category. The map $R:D\to C$ sends every
old-labelled arrow to its label and $1_{x'}$ to $1_x$; it is
a functor. The inclusion $I:C\to D$ is a functor and $RI=1_C$.
The old identity label gives an arrow $u:x\to x'$ in $D$.
With component $u$ at $x'$ and identity components elsewhere,
we have a natural transformation $IR\Rightarrow1_D$.
Naturality is exactly the identity law in $C$ for old-labelled
arrows, and the new identity causes no additional equation.
Thus $BI$ and $BR$ are homotopy inverses.

Let $v:x'\to x$ be the other typed copy of $1_x$. Then
$vu=1_x$, whereas $uv$ is the old identity label in
$D(x',x')$. Hence $x$ is a retract of $x'$. The full
subcategory $D\setminus\{x\}$ has the same classifying
space as $D$ by Lemma~\ref{CP2:A-retract-deletion}; after
renaming $x'$ as $x$, it is exactly $C^+$.
This proves the claimed homotopy equivalence and all count
assertions. Iterating at the chosen corner gives every larger
arrow count. No inclusion $C\to C^+$ that sends old arrows
to themselves is asserted to preserve identities.
\end{proof}
